\documentclass[11pt,letterpaper]{article}
\usepackage[margin=0.85in,headheight=14pt]{geometry}
\usepackage{amsmath,amssymb}
\usepackage{enumitem}
\usepackage{fancyhdr}
\usepackage{lmodern}
\usepackage[T1]{fontenc}
\usepackage[dvipsnames]{xcolor}
\usepackage[most]{tcolorbox}

\pagestyle{fancy}
\fancyhf{}
\lhead{SYDE 556/750 \qquad Recipe practice}
\rhead{Solutions}
\cfoot{\thepage}
\renewcommand{\headrulewidth}{0pt}

\setlength{\parindent}{0pt}
\setlength{\parskip}{4pt}

\newcommand{\Jb}{J^{\mathrm{bias}}}
\newcommand{\ans}[1]{\par\vspace{2pt}{\bfseries #1}\quad}

\newtcolorbox{result}{colback=Green!6,colframe=Green!50!black,boxrule=0.5pt,left=4pt,right=4pt,top=2pt,bottom=2pt,arc=1pt}

\setlist[enumerate]{leftmargin=1.5em,itemsep=2pt,topsep=2pt}

\begin{document}

\begin{center}
{\Large\bfseries Recipe practice solutions}\\[2pt]
{\large Feedforward transformations}\\[4pt]
Open a solution only after attempting that question. Each green box is the result.
\end{center}

%---------------------------------------------------------------------
\section{One pass through all four steps}

\ans{1.}
The input currents are linear in the represented scalar:
\[
J_1^{A} = x + 1, \qquad J_2^{A} = -x + 1.
\]
The output currents are the same kind of encoding, written in $y$ rather than $x$:
\[
J_1^{B} = 2y + \tfrac12, \qquad J_2^{B} = -y + 1.
\]
The activities of A are already given, so these input currents are not turned back into rates.
\begin{result}
$J_1^A = x + 1$, $J_2^A = -x + 1$. \quad $J_1^B = 2y + 1/2$, $J_2^B = -y + 1$.
\end{result}

\ans{2.}
The value represented by B is twice the value represented by A:
\[
y = 2x.
\]
\begin{result}
$y = 2x$.
\end{result}

\ans{3.}
Rows are neurons and columns are the inputs $x = 0$ and $x = 1$:
\[
A = \begin{pmatrix} 1 & 2 \\ 2 & 1 \end{pmatrix}, \qquad
Y = \begin{pmatrix} 0 & 2 \end{pmatrix}.
\]
\[
AA^{\mathsf T} = \begin{pmatrix} 5 & 4 \\ 4 & 5 \end{pmatrix}, \qquad
(AA^{\mathsf T})^{-1} = \frac19 \begin{pmatrix} 5 & -4 \\ -4 & 5 \end{pmatrix}, \qquad
AY^{\mathsf T} = \begin{pmatrix} 4 \\ 2 \end{pmatrix}.
\]
\[
(D^f)^{\mathsf T}
  = \frac19 \begin{pmatrix} 5 & -4 \\ -4 & 5 \end{pmatrix}
    \begin{pmatrix} 4 \\ 2 \end{pmatrix}
  = \begin{pmatrix} 4/3 \\ -2/3 \end{pmatrix},
\qquad
D^f = \begin{pmatrix} 4/3 & -2/3 \end{pmatrix}.
\]
On the two samples, $a = (1, 2)$ and $a = (2, 1)$:
\[
D^f a = 0 \quad \text{and} \quad D^f a = 2,
\]
which are $2x$ at $x = 0$ and $x = 1$. Two neurons and two samples, with this invertible $A$, can hit both targets exactly.
\begin{result}
$D^f = (4/3,\ -2/3)$. It reproduces $(0,\ 2)$.
\end{result}

\ans{4.}
The output encoders, with gains included, are the rows of
\[
E = \begin{pmatrix} 2 \\ -1 \end{pmatrix}.
\]
\[
W^f = E D^f
  = \begin{pmatrix} 2 \\ -1 \end{pmatrix}
    \begin{pmatrix} 4/3 & -2/3 \end{pmatrix}
  = \begin{pmatrix} 8/3 & -4/3 \\ -4/3 & 2/3 \end{pmatrix}.
\]
Biases stay outside the matrix:
\[
J_1^{B} = \tfrac83 a_1 - \tfrac43 a_2 + \tfrac12, \qquad
J_2^{B} = -\tfrac43 a_1 + \tfrac23 a_2 + 1.
\]
At $x = 1$ the activities are $a_1 = 2$, $a_2 = 1$, and $2x = 2$.
From the weights,
\[
J_1^{B} = \tfrac83\cdot 2 - \tfrac43\cdot 1 + \tfrac12 = \tfrac92, \qquad
J_2^{B} = -\tfrac43\cdot 2 + \tfrac23\cdot 1 + 1 = -1.
\]
From the output encoding with $y$ replaced by $2x$,
\[
J_1^{B} = 2\cdot 2 + \tfrac12 = \tfrac92, \qquad
J_2^{B} = -2 + 1 = -1.
\]
Same currents. That is the substitution the fourth step asks for.
\begin{result}
$W^f = \begin{pmatrix} 8/3 & -4/3 \\ -4/3 & 2/3 \end{pmatrix}$.
At $x = 1$ both routes give $J = (9/2,\ -1)$.
\end{result}

%---------------------------------------------------------------------
\section{Same populations, new function}

\ans{1.}
The encodings of A and of B were fixed by the populations, so step 1 is unchanged. Step 2 is a new sentence, $y = x + 1$. Steps 3 and 4 use that new target, so $Y$, $D^f$, and $W^f$ are recomputed. $E$ and the biases belong to B and stay put.
\begin{result}
Copy step 1, and copy $E$. Redo the transformation, $D^f$, and $W^f$.
\end{result}

\ans{2.}
\[
Y = \begin{pmatrix} 1 & 2 \end{pmatrix}, \qquad
AY^{\mathsf T} = \begin{pmatrix} 5 \\ 4 \end{pmatrix}.
\]
\[
(D^f)^{\mathsf T}
  = \frac19 \begin{pmatrix} 5 & -4 \\ -4 & 5 \end{pmatrix}
    \begin{pmatrix} 5 \\ 4 \end{pmatrix}
  = \begin{pmatrix} 1 \\ 0 \end{pmatrix},
\qquad
D^f = \begin{pmatrix} 1 & 0 \end{pmatrix}.
\]
The second input neuron receives a decoder of zero for this particular target. That is a fact about these samples, not a rule about constant offsets.
\begin{result}
$D^f = (1,\ 0)$.
\end{result}

\ans{3.}
\[
W^f = \begin{pmatrix} 2 \\ -1 \end{pmatrix}
      \begin{pmatrix} 1 & 0 \end{pmatrix}
    = \begin{pmatrix} 2 & 0 \\ -1 & 0 \end{pmatrix}.
\]
\[
J_1^{B} = 2 a_1 + \tfrac12, \qquad
J_2^{B} = -a_1 + 1.
\]
At $x = 1$, $a_1 = 2$ and $y = 2$, so $J_1^{B} = 9/2$ and $J_2^{B} = -1$, which is $\alpha e y + \Jb$ for each neuron of B.
\begin{result}
$J_1^B = 2a_1 + 1/2$, \ $J_2^B = -a_1 + 1$.
\end{result}

\ans{4.}
\[
d + d^{\mathbf 1}
  = \bigl(\tfrac23,\ -\tfrac13\bigr) + \bigl(\tfrac13,\ \tfrac13\bigr)
  = (1,\ 0),
\]
which is the decoder from part 2. Least squares is linear in the target, and $x + 1$ is the identity plus the constant function. The identity decoder alone reconstructs $x$. Substituting that into B would encode $x$, and the $+1$ would be missing.
\begin{result}
$d + d^{\mathbf 1} = (1,\ 0)$. The identity decoder implements $y = x$.
\end{result}

\end{document}
